\section{Discussion}
\label{sec:discussion}

\para{What the results show.} The findings support all four sub-hypotheses that
motivated the framework. The hybrid beats every baseline with non-overlapping
bootstrap CIs and significant Holm-corrected McNemar tests (H1). The two single
modalities fail on disjoint, complementary slices---the Bayesian view on design
flaws (H2), the symbolic view on power (H3)---and fusion wins precisely because it
is the only method that observes all the discriminating inputs at once: sample
size and noise from the raw observations, and structural flags and intervention
strength from the experiment's logic. This is the central epistemic claim of the
work, made measurable: an evidence layer and a structure layer are each necessary
but, alone, insufficient. The framework also produces interpretable, calibrated
outputs (H4): every diagnosis is a short reasoning chain that names the operative
check, and confidence is monotone, if conservative.

\para{Does the structure actually do work?} The most informative comparison is the
frontier LLM judge. Given identical information and an explicit rubric, GPT-4.1
reaches $0.888$---it nails design flaws ($0.97$, the flags are textually salient)
but still confuses robustness with weak intervention because it does not reliably
translate a small $N$ or a noisy metric into ``underpowered''. The explicit
power/identifiability logic of the framework supplies exactly the inference a
strong general reasoner does not provide for free, and does so transparently and
at near-zero marginal cost. The negative result for \nhst{} (accuracy $0.36$) is
the empirical face of ``absence of evidence is not evidence of
absence''~\citep{dale2024confirming}.

\para{Connection to the literature.} The design operationalises the reviewed
pillars: equivalence-margin and confirmability gating from
\citet{dale2024confirming}, identifiability as an equivalence-class obstruction
from \citet{santosgrueiro2026limits}, Bayes-factor evidence with prior sensitivity
from \citet{rouder2009bayesian, robert2016demise}, and graded red-teaming as the
intervention-strength axis anchored on \harmbench~\citep{mazeika2024harmbench}.

\para{Limitations.} We are explicit about the scope of these numbers.
\begin{itemize}[leftmargin=*,itemsep=1pt,topsep=1pt]
  \item \textbf{Synthetic benchmark, separable by construction.} The generator
        encodes the discriminating structure that fusion exploits, so the absolute
        $0.99$ will not transfer verbatim to messy real experiments. This is a
        controlled proof-of-mechanism that demonstrates and quantifies the fusion
        logic and each modality's blind spot---not a field accuracy estimate. The
        literature-confirmed absence of a cause-labelled real benchmark is exactly
        the gap that motivates building one.
  \item \textbf{Dataset-informed thresholds.} The adequacy checks
        ($N \ge 50$, residual $\mathrm{sd} \le 2.0$, strength $\ge 0.5$) are
        interpretable rules of thumb tuned to the generator; on real data they
        would be set by a proper power analysis and measurement-validity audit per
        study. We show robustness to the \emph{prior} (\secref{sec:sensitivity});
        robustness to \emph{these thresholds} is future work.
  \item \textbf{Single LLM judge, no human panel.} GPT-4.1 is a proxy for the
        intended human expert; true inter-annotator $\kappa$ against human
        methodologists is not measured here.
  \item \textbf{Under-confidence.} The hybrid's ECE of $0.248$ is conservative;
        confidences would benefit from post-hoc calibration (\eg isotonic
        regression) before being surfaced to users as probabilities.
  \item \textbf{Hard three-cause taxonomy.} Real nulls can have mixed or compound
        causes; the framework currently emits a single hard label (plus reasoning),
        not a multi-label or partial-cause decomposition.
\end{itemize}

\para{Broader implications.} A diagnostic that says \emph{which} of the three a
null is---with a reason---directly improves the reliability of alignment science.
It guards against the most dangerous failure mode, shipping false confidence in a
model's safety when a red-team attack was merely too weak or a confound masked a
real vulnerability, and it turns each null into an actionable next step rather than
an ambiguous stopping point. Because the framework is interpretable and cheap, it
can sit inside an existing evaluation workflow as a sanity check on every
no-effect finding.
